\subsection{Pregunta 1}
¿Qué peligros puede haber en una red pública y nmap?

\begin{respuesta}
En una red pública cualquier persona conectada puede interceptar el tráfico de otros usuarios ya que no existe cifrado garantizado entre los dispositivos. Esto permite ataques como el man-in-the-middle, donde un atacante se pone entre dos equipos y lee la comunicación. Usar nmap en una red pública sin autorización puede ser ilegal, ya que escanear puertos de dispositivos ajenos equivale a explorar su infraestructura sin permiso. Además un atacante podría usar nmap para descubrir qué dispositivos tienen puertos vulnerables y lanzar ataques dirigidos contra ellos.


    
\end{respuesta}
